\documentclass[xcolor=dvipsnames,professionalfonts]{beamer}


\usepackage{xcolor}
\usefonttheme{serif}

\usepackage{tikz}


\usepackage{xepersian}
\settextfont{Yas.ttf}

\usepackage{mytheme}

\usebackgroundtemplate{\includegraphics[width=\paperwidth, height=\paperheight]{background-brightness-3.jpg}}

\author{علی ابراهیمیان}
\title{پرزنتیشن}
\institute{دانشگاه آزاد اسلامی}

\newtheorem{thm}{قضیه}

\begin{document}
    \section{بخش ۱}
    \frame{\maketitle}
    \begin{frame}
        \frametitle{عنوان اسلاید اول}
        \begin{block}

            
        این اولین صفحه اسلایدر است.
    
            \begin{itemize}\raggedright
            \item گزینه اول
            \item گزینه دوم
            \item گزینه سوم
            \end{itemize}
        \end{block}
    \end{frame}

    \begin{frame}
        \frametitle{عنوان اسلاید دوم}
        \begin{block}

            این یک فرمول است

            \[
            \int_{0}^{1}{f(x)}dx
            \]
        \end{block}
    \end{frame}

    \begin{frame}
        \frametitle{وقفه برای رفتن به  متن بعدی}
        \begin{block}{تستی}
            این متن تستی است
        \begin{itemize}
            \item علی
            \item مهدی
            \pause
            \item هادی
            \pause
            
        \end{itemize}
        \end{block}
    \end{frame}

    \begin{frame}
        \frametitle{ریاضی}
        \begin{thm}
        فرض کنید
        \begin{itemize}\raggedright
            \item تابع مشتق پذیر باشد.
        \end{itemize}
        \end{thm}
    \end{frame}
\end{document}